\chapter{Assignment 3 --- Low-Pass Filtering of an Image}
\label{ch:hw3}

\section{The problem}
\label{sec:hw3-problem}

A low-pass filter smooths an image. The provided image is smoothed with three
box filters of size $3\times3$, $5\times5$ and $7\times7$, implemented in the
\texttt{conv} function of the provided template
\texttt{Low\_pass\_filter\_template.py}. New code was added, but the template
itself was not changed.

A box filter replaces the value at the centre of the box with the average of the
values inside it. For the $3\times3$ case the filter matrix is

\begin{equation}
\label{eq:hw3-filter3}
\mathit{filter} = \frac{1}{9}
\begin{bmatrix}
1 & 1 & 1 \\
1 & 1 & 1 \\
1 & 1 & 1
\end{bmatrix}
=
\begin{bmatrix}
\frac{1}{9} & \frac{1}{9} & \frac{1}{9} \\[2pt]
\frac{1}{9} & \frac{1}{9} & \frac{1}{9} \\[2pt]
\frac{1}{9} & \frac{1}{9} & \frac{1}{9}
\end{bmatrix}.
\end{equation}

The assignment's worked example takes the sub-image

\begin{equation}
\label{eq:hw3-subimage}
\mathit{subImage}_1 =
\begin{bmatrix}
21 & 78 & 63 \\
16 & 109 & 73 \\
6 & 88 & 86
\end{bmatrix},
\end{equation}

and multiplies it element-wise by the filter,

\begin{equation}
\label{eq:hw3-conv}
C = \mathit{filter} * \mathit{subImage}_1 =
\begin{bmatrix}
2.33 & 8.67 & 7 \\
1.78 & 12.11 & 8.11 \\
0.67 & 9.78 & 9.56
\end{bmatrix}.
\end{equation}

Summing the entries of $C$ should reproduce the mean of $\mathit{subImage}_1$,
because that is exactly what a box filter is defined to compute. The check is
done one operation at a time in \Cref{eq:hw3-checksum}.

\begin{align*}
&2.33 + 8.67 + 7      = 18   && \text{row 1 of } C \\
&1.78 + 12.11 + 8.11  = 22   && \text{row 2 of } C \\
&0.67 + 9.78 + 9.56   = 20   && \text{row 3 of } C \\
&18 + 22 + 20         = 60   && \text{sum of } C
\end{align*}

\begin{align}
\label{eq:hw3-checksum}
&21+78+63+16+109+73+6+88+86 = 540 && \text{sum of } \mathit{subImage}_1 \\
&540 / 9 = 60 && \text{mean of } \mathit{subImage}_1 \notag
\end{align}

The two routes agree at $60$, confirming that the filter matrix does perform
the averaging. Running the finished \texttt{conv} function
(\Cref{sec:hw3-impl}) on $\mathit{subImage}_1$ with the $3\times3$ filter
returns exactly $60.0$, matching this hand calculation. The $5\times5$ and
$7\times7$ filters are built the same way, with $1/25$ and $1/49$ in every
entry.

\section{Implementation}
\label{sec:hw3-impl}

The template left two gaps: \texttt{lowPassFilter}, which builds the filter
matrix, and \texttt{conv}, which slides that filter over the image.

\paragraph{Problem: the filter entries must sum to 1.} If the filter's entries
summed to anything other than $1$, every output pixel would be a value scaled
away from the true local average, and the whole image would come out
uniformly brighter or darker regardless of what it originally showed. An
averaging filter must leave the image's overall brightness unchanged; only the
high-frequency detail (fine texture, grain) should be removed. \texttt{lowPassFilter} solves this by filling an $\mathit{lpfw}\times\mathit{lpfh}$
matrix with the single value $1/(\mathit{lpfw}\cdot\mathit{lpfh})$ in every
entry:

\begin{equation}
\label{eq:hw3-lowpassfilter}
\texttt{lowPass = np.ones((lpfw, lpfh)) / (lpfw * lpfh)}.
\end{equation}

Every entry is $1/(\mathit{lpfw}\cdot\mathit{lpfh})$ and there are
$\mathit{lpfw}\cdot\mathit{lpfh}$ of them, so the entries sum to exactly $1$
by construction, for any box size.

\paragraph{Problem: what happens at the boundary of the image.} A filter of
size $\mathit{fw}\times\mathit{fh}$ needs a full $\mathit{fw}\times\mathit{fh}$
neighbourhood centred at each output pixel. Near the edge of the image that
neighbourhood runs off the array, so there is no well-defined average to
compute there. \texttt{conv} solves this by simply not producing an output for
those positions: the sliding window only visits positions where the full
window still fits inside the input, which is why the output array is declared
smaller than the input,

\begin{equation}
\label{eq:hw3-outshape}
\mathit{out.shape} = (\mathit{iw} - \mathit{fw} + 1,\; \mathit{ih} - \mathit{fh} + 1),
\end{equation}

and the double loop in \texttt{conv} fills each position with the element-wise
sum of the filter against the matching sub-image:

\begin{equation}
\label{eq:hw3-convloop}
\texttt{out[i, j] = np.sum(image[i:i+fw, j:j+fh] * filter)}.
\end{equation}

For the $512\times512$ input this drops the boundary pixels, giving
$510\times510$, $508\times508$ and $506\times506$ outputs for the $3\times3$,
$5\times5$ and $7\times7$ filters respectively (\Cref{tab:hw3-results}).

The only other addition beyond \texttt{lowPassFilter} and \texttt{conv} is a
guard placed right after the image is read,

\begin{equation}
\label{eq:hw3-guard}
\texttt{if img.ndim == 3: img = img.mean(axis=2)},
\end{equation}

because the template's \texttt{iw, ih = image.shape} assumes a single-channel
image and would fail on a 3-channel one. This is the one place where the
template's assumption is fragile: for this input the guard did not actually
trigger, since \texttt{original.jpg} is a baseline single-component JPEG and
\texttt{mpimg.imread} already returned a $(512, 512)$ uint8 array, but the
guard protects against inputs that are not.

No normalization step was needed either. The image is a uint8 array and the
filter is a float64 array, so their product upcasts to float64; because the
filter entries sum to $1$ (\Cref{eq:hw3-lowpassfilter}), the averaged values
never leave the range $0$ to $255$, and \texttt{cv2.imwrite} saturate-casts the
result back to 8 bits when writing the file (it prints a benign
``Unsupported depth \ldots fallbacked to CV\_8U'' warning in the process, shown
in \Cref{lst:hw3-out}).

The complete program is listed below; it is the same file submitted for the
assignment.

\lstinputlisting[style=py,
                 caption={Lim\_20701313\_hw3.py: the complete solution.},
                 label={lst:hw3-code}]{../hw3/submission/Lim_20701313_hw3.py}

\section{Results}
\label{sec:hw3-results}

\Cref{fig:hw3-outputs} shows the original image next to the three filtered
outputs. \Cref{tab:hw3-results} lists the corresponding filter entry value and
output shape for each box size.

\begin{figure}[ht]
\centering
\begin{minipage}{0.48\linewidth}
\centering
\includegraphics[width=\linewidth]{hw3_original.png}
\par\textbf{Original} ($512\times512$).
\end{minipage}
\hfill
\begin{minipage}{0.48\linewidth}
\centering
\includegraphics[width=\linewidth]{hw3_low3x3.png}
\par\textbf{$3\times3$ box filter} ($510\times510$).
\end{minipage}

\vspace{1em}

\begin{minipage}{0.48\linewidth}
\centering
\includegraphics[width=\linewidth]{hw3_low5x5.png}
\par\textbf{$5\times5$ box filter} ($508\times508$).
\end{minipage}
\hfill
\begin{minipage}{0.48\linewidth}
\centering
\includegraphics[width=\linewidth]{hw3_low7x7.png}
\par\textbf{$7\times7$ box filter} ($506\times506$).
\end{minipage}
\caption{Original image and low-pass filtered outputs at three box sizes.}
\label{fig:hw3-outputs}
\end{figure}

\begin{table}[ht]
\centering
\caption{Filter size, entry value and output shape for each run.}
\label{tab:hw3-results}
\begin{tabular}{lll}
\toprule
Box size & Filter entry value & Output shape \\
\midrule
$3\times3$ & $1/9 \approx 0.11111111$  & $510\times510$ \\
$5\times5$ & $1/25 = 0.04$             & $508\times508$ \\
$7\times7$ & $1/49 \approx 0.02040816$ & $506\times506$ \\
\bottomrule
\end{tabular}
\end{table}

The console output of the run is reproduced verbatim, with the $5\times5$ and
$7\times7$ filter matrices abbreviated since every entry is identical:

\begin{lstlisting}[style=shell,caption={Console output of the program.},label={lst:hw3-out}]
(512, 512)
 ================ Low Pass Filter ===================
Please create a (lpfw x lpfh) Low Pass Filter matrix
[[0.11111111 0.11111111 0.11111111]
 [0.11111111 0.11111111 0.11111111]
 [0.11111111 0.11111111 0.11111111]]
Please create a (lpfw x lpfh) Low Pass Filter matrix
[[0.04 0.04 ... ]]   # 5x5, every entry 0.04
Please create a (lpfw x lpfh) Low Pass Filter matrix
[[0.02040816 ... ]]  # 7x7, every entry 0.02040816
\end{lstlisting}

The smoothing grows with the box size, exactly as expected from averaging over
more neighbours per output pixel. At $3\times3$ the film grain is only lightly
softened and fine detail such as hair strands, eyelashes and skin texture
stays sharp. At $5\times5$ the grain is largely gone and fine hair strands
start merging, while the eyes, nose and lips stay well defined. At $7\times7$
the smoothing is strongest: the background and skin flatten, hair strands
merge into soft masses, and edges such as the jawline show visible softening.
